% GENERATED FILE. Edit canonical lessons, then run npm run generate:books.
% canonical-source-hash: fnv1a64:77fe3df817718948
% canonical-lessons: KA-C110-sante, KA-C110-kadalatira, KA-C110-tota, KA-C110-mule, KA-C110-tiruvu

\chapter{Market, Beach, Garden, Corner, Turn}
\label{ch:ka-market-beach-garden-corner-turn}
\hlchaptermodality{110}

\begin{chapteropening}
\textbf{By the end of this chapter:} \emph{I can say a market, a beach, a garden, a corner and a turn.}

Complete the last lesson of chapter 110: i can say a market, a beach, a garden, a corner and a turn.
\end{chapteropening}

\section[sante]{\kn{ಸಂತೆ} (sante) — a market}

\label{lesson:KA-C110-sante}



\begin{quote}
Before the new one: say the Kannada for a library, then the Kannada for a bank.
\end{quote}

\subsection*{You'll want to know: \kn{ಸಂತೆ}}
\textbf{\kn{ಸಂತೆ}} — \emph{sante} — \textquotedblleft{}a market\textquotedblright{}.

\begin{grammarlens}[title={What you've built}]
One more word for this chapter.
\end{grammarlens}

\subsection*{Your turn}
\begin{itemize}
\raggedright
  \item \emph{Say it:} \emph{sante}
  \item \emph{Say it:} \emph{sante}, once more
  \item \emph{Say it:} say \emph{sante}
  \item \emph{Recall:} say the Kannada for a library, then the Kannada for a bank, then say \emph{sante} again
\end{itemize}

\begin{tcolorbox}[breakable,colback=teal!4,colframe=teal!35!black,title={Before you move on}]
What does \kn{ಸಂತೆ} mean? (\textquotedblleft{}A market\textquotedblright{}.) Say it once more.
\end{tcolorbox}
\section[kaḍalatīra]{\kn{ಕಡಲತೀರ} (kaḍalatīra) — a beach}

\label{lesson:KA-C110-kadalatira}



\begin{quote}
Before the new one: say the Kannada for a bank, then the Kannada for a market.
\end{quote}

\subsection*{You'll want to know: \kn{ಕಡಲತೀರ}}
\textbf{\kn{ಕಡಲತೀರ}} — \emph{kaḍalatīra} — \textquotedblleft{}a beach\textquotedblright{}.

\begin{grammarlens}[title={What you've built}]
One more word for this chapter.
\end{grammarlens}

\subsection*{Your turn}
\begin{itemize}
\raggedright
  \item \emph{Say it:} \emph{kaḍalatīra}
  \item \emph{Say it:} \emph{kaḍalatīra}, once more
  \item \emph{Say it:} say \emph{kaḍalatīra}
  \item \emph{Recall:} say the Kannada for a bank, then the Kannada for a market, then say \emph{kaḍalatīra} again
\end{itemize}

\begin{tcolorbox}[breakable,colback=teal!4,colframe=teal!35!black,title={Before you move on}]
What does \kn{ಕಡಲತೀರ} mean? (\textquotedblleft{}A beach\textquotedblright{}.) Say it once more.
\end{tcolorbox}
\section[tōṭa]{\kn{ತೋಟ} (tōṭa) — a garden}

\label{lesson:KA-C110-tota}



\begin{quote}
Before the new one: say the Kannada for a market, then the Kannada for a beach.
\end{quote}

\subsection*{You'll want to know: \kn{ತೋಟ}}
\textbf{\kn{ತೋಟ}} — \emph{tōṭa} — \textquotedblleft{}a garden\textquotedblright{}.

\begin{grammarlens}[title={What you've built}]
One more word for this chapter.
\end{grammarlens}

\subsection*{Your turn}
\begin{itemize}
\raggedright
  \item \emph{Say it:} \emph{tōṭa}
  \item \emph{Say it:} \emph{tōṭa}, once more
  \item \emph{Say it:} say \emph{tōṭa}
  \item \emph{Recall:} say the Kannada for a market, then the Kannada for a beach, then say \emph{tōṭa} again
\end{itemize}

\begin{tcolorbox}[breakable,colback=teal!4,colframe=teal!35!black,title={Before you move on}]
What does \kn{ತೋಟ} mean? (\textquotedblleft{}A garden\textquotedblright{}.) Say it once more.
\end{tcolorbox}
\section[mūle]{\kn{ಮೂಲೆ} (mūle) — a corner}

\label{lesson:KA-C110-mule}



\begin{quote}
Before the new one: say the Kannada for a beach, then the Kannada for a garden.
\end{quote}

\subsection*{You'll want to know: \kn{ಮೂಲೆ}}
\textbf{\kn{ಮೂಲೆ}} — \emph{mūle} — \textquotedblleft{}a corner\textquotedblright{}.

\begin{grammarlens}[title={What you've built}]
One more word for this chapter.
\end{grammarlens}

\subsection*{Your turn}
\begin{itemize}
\raggedright
  \item \emph{Say it:} \emph{mūle}
  \item \emph{Say it:} \emph{mūle}, once more
  \item \emph{Say it:} say \emph{mūle}
  \item \emph{Recall:} say the Kannada for a beach, then the Kannada for a garden, then say \emph{mūle} again
\end{itemize}

\begin{tcolorbox}[breakable,colback=teal!4,colframe=teal!35!black,title={Before you move on}]
What does \kn{ಮೂಲೆ} mean? (\textquotedblleft{}A corner\textquotedblright{}.) Say it once more.
\end{tcolorbox}
\section[tiruvu]{\kn{ತಿರುವು} (tiruvu) — a turn, a bend}

\label{lesson:KA-C110-tiruvu}



\begin{quote}
Before the new one: say the Kannada for a garden, then the Kannada for a corner.
\end{quote}

\subsection*{You'll want to know: \kn{ತಿರುವು}}
\textbf{\kn{ತಿರುವು}} — \emph{tiruvu} — \textquotedblleft{}a turn, a bend\textquotedblright{}.

\begin{grammarlens}[title={What you've built}]
One more word for this chapter.
\end{grammarlens}

\subsection*{Your turn}
\begin{itemize}
\raggedright
  \item \emph{Say it:} \emph{tiruvu}
  \item \emph{Say it:} \emph{tiruvu}, once more
  \item \emph{Say it:} say \emph{tiruvu}
  \item \emph{Recall:} say the Kannada for a garden, then the Kannada for a corner, then say \emph{tiruvu} again
\end{itemize}

\begin{tcolorbox}[breakable,colback=teal!4,colframe=teal!35!black,title={Before you move on}]
What does \kn{ತಿರುವು} mean? (\textquotedblleft{}A turn, a bend\textquotedblright{}.) Say it once more.
\end{tcolorbox}
